\section{Analisis Celah Kebijakan (Policy Gap)}
\label{sec:celah}

\subsection{Kerangka Regulasi yang Berlaku}

Tabel~\ref{tab:regulasi} membandingkan regulasi dan rencana yang berlaku dengan kebutuhan solusi terpilih.

\begin{table}[H]
\centering
\caption{Regulasi terkait dan celahnya terhadap solusi integrasi}
\label{tab:regulasi}
\small
\begin{tabularx}{\textwidth}{L{3.6cm} Y Y}
\toprule
Regulasi atau rencana & Yang diatur & Celah terhadap solusi integrasi \\
\midrule
Perpres 112/2022 \cite{R-PERPRES-112-2022} & Percepatan energi terbarukan, larangan PLTU baru, peta jalan pengakhiran PLTU & Peta jalan tidak menyebut sumber energi pengganti per PLTU \\
\addlinespace
Permen ESDM 10/2025 \cite{R-PERMEN-ESDM-10-2025} & Kriteria pensiun dini Pasal 11 dan kewajiban kajian oleh PLN & Kriteria belum menimbang jenis kontrak (IPP atau PLN) dan dampak kesehatan per lokasi. Hasil kajian PLN tidak dibuka \\
\addlinespace
Permen ESDM 2/2024 \cite{R-PERMEN-ESDM-2-2024} & PLTS atap terhubung jaringan, kuota per wilayah, ekspor tidak mengurangi tagihan & Kuota ditetapkan tanpa dikaitkan dengan PLTU yang akan dipensiunkan \\
\addlinespace
Kuota PLTS atap hingga 2028 \cite{B-CNBC-ROOFTOP-2028} & 2 GW nasional, 1,85 GW Jawa-Madura-Bali & Realisasi tertinggal, belum ada kewajiban bagi gedung pemerintah dan BUMN \\
\addlinespace
CIPP JETP \cite{P-JETP-CIPP-2023} & Target pensiun dini 1,7 GW hingga 2040 & Alokasi dana kecil dan tidak dipersyaratkan terikat sebelum penutupan \\
\addlinespace
RUPTL 2025-2034 \cite{B-LISTRIKID-RUPTL-JAMALI} & Tambahan pembangkit Jamali termasuk surya skala besar & Pembangkit pengganti tidak dijadwalkan terhadap tanggal penutupan PLTU tertentu \\
\bottomrule
\end{tabularx}
\end{table}

\subsection{Celah Kebijakan Kritis}

Berdasarkan perbandingan di atas, kajian ini mengidentifikasi enam celah yang saling berkaitan.

\begin{enumerate}
  \item \textbf{Celah keterpaduan.} PLTS atap dan pensiun dini diatur oleh dua peraturan yang tidak saling merujuk. Akibatnya, PLTS atap dihitung PLN sebagai pengurang pendapatan, bukan sebagai pengganti PLTU yang membuka penghematan.
  \item \textbf{Celah pemilihan PLTU.} Pensiun dini PLTU IPP menghentikan pembayaran kapasitas, sedangkan pensiun dini PLTU milik PLN tidak. Perbedaan ini menentukan apakah pensiun dini menguntungkan atau membebani negara (Bagian~\ref{sec:kelayakan}), tetapi belum menjadi pertimbangan eksplisit dalam kriteria.
  \item \textbf{Celah pendanaan.} Belum ada mekanisme yang memastikan kompensasi dari dana luar terikat sebelum keputusan penutupan. Pembatalan Cirebon-1 terjadi pada tahap ini.
  \item \textbf{Celah keandalan.} Reserve margin Jawa-Bali diperkirakan turun ke batas aman 33\% sekitar 2027-2029 \cite{B-PETROMINDO-RESERVE}. Setelah itu, setiap PLTU yang pensiun memerlukan pembangkit pengganti yang siap lebih dulu. PLTS atap tidak menambah daya mampu saat beban puncak malam hari.
  \item \textbf{Celah transisi berkeadilan.} Program pelatihan ulang belum dijadwalkan terhadap tanggal penutupan setiap PLTU.
  \item \textbf{Celah data dan transparansi.} Kajian pensiun dini PLN dan kontrak IPP tidak terbuka. Pemangku kepentingan tidak memiliki alat bersama untuk menguji siapa menanggung biaya dan manfaat suatu usulan, sehingga perbedaan penilaian seperti pada Cirebon-1 sulit diselesaikan.
\end{enumerate}

Keenam celah ini menunjukkan bahwa persoalannya bukan kekurangan teknologi. Yang belum ada adalah aturan yang memasangkan PLTS atap dengan pensiun dini, urutan pemilihan PLTU yang menguntungkan negara, syarat penutupan yang jelas, dan alat hitung bersama yang dapat dipercaya semua pihak.
